\begin{frame}
\frametitle{Conclusioni finali}
In queste diapositive si è quindi presentato lo sviluppo e l'implementazione 
di un canale nascosto tramite il protocollo ICMP. L’obbiettivo era dimostrare 
come il protocollo, fondamentale per il funzionamento delle reti, potesse 
essere sfruttabile per la creazione di un Covert Channel che possa esfiltrare 
delle informazioni. 
\vspace{1ex} \newline 
Si è descritto l’implementazione della comunicazione fra le entità presenti 
nel sistema e successivamente l'implementazione delle seguenti tipologie di 
Cover Channel: 
\begin{itemize}
    \item Storage Covert Channel: inserendo i dati ne campi dei 
    messaggi ICMP. 
    \item Timing Covert Channel: codificando le informazioni 
    tramite una funzione 
    %$\text{min\_delay}+(\text{byte}/255)*(\text{max\_delay}-\text{min\_delay})$
    al cui valore finale verrà agiunto del rumore casuale.  
    \item Behavioural Covert Channel: codifica i dati attraverso 
    degli eventi definiti dalla coppia Tipologia e Codice del 
    messaggio ICMP.  
\end{itemize} 
\end{frame} 
\begin{frame}
Dopodiché si sono testati i Covert Channel tramite RITA e dalle analisi 
risulta che la rilevabilità di un Covert channel dipende dai seguenti fattori: 
\begin{itemize}
    \item %L'utilizzo di un delay fra l'invio di un blocco di dati ed il 
    %successivo permette di abbassare la visibilità del Covert Channel. 
    La frequenza con cui si inviano i dati. 
    \item La quantità di pacchetti mandati da una sorgente in uno o più 
    intervalli di tempo. 
\end{itemize} 
Sebbene sia un protocollo funzionale al monitoraggio di rete, i messaggi 
ICMP devono essere limitati solo a quelle tipologie essenziali alla rete 
E fra queste bisogna monitorare il loro utilizzo ed il loro 
contenuto. Ciò richiederà un ispezione dei pacchetti presenti nel flusso, 
un analisi statistica per la ricerca di pattern e la limitazione nella 
rete dei messaggi ICMP. 
\end{frame} 

